\section{Ethics Considerations}\label{sec:ethics}


\noindent \textbf{Controlled sandbox execution.} All model loading and inference, including the malicious corpus, is performed in an isolated sandbox on our own workstation. No production system, third-party service, or end-user device was targeted, and no real credentials, tokens, or private data were exposed during evaluation. The network endpoints contacted by the synthetic payloads are local and under our control, so no external party is affected.

\noindent \textbf{Legitimate use of public model hubs.} The 201 benign Part~A models are publicly available artifacts downloaded from the Hugging Face Hub strictly for false-positive profiling; we neither modify nor redistribute them, and we respect their original licenses. Platform-hosted scanner comparisons (JFrog, Protect~AI, ClamAV) were conducted by uploading our own artifacts through Hugging Face's intended public scanning pipeline, which is the sanctioned use of the platform's security API rather than an attack on the platform itself.

\noindent \textbf{Dual-use and responsible release.} As with any detection system, \toolname{} and the accompanying benchmark are dual-use artifacts: the attack chains we describe could, in principle, inform adversarial efforts. To mitigate this risk, we publicly release all models in our constructed datasets, but every payload is neutered so that it executes only harmless commands and is incapable of real exfiltration, persistence, or system compromise. We strongly recommend that anyone reproducing or extending our work execute these models exclusively within a controlled sandbox environment isolated from production systems, sensitive credentials, and external networks, so that unintended behavior cannot propagate beyond the test host. The benchmark is intended solely for defensive security research.
